\DO{It is common to use subsections for the separate pieces that
  together make up the section.}

When you describe the details of your work, you might sometimes need
to include an equation. If you, for example, want to describe the
Pythagorean Theorem you would state
\begin{align}
  \label{eqn:pythagoras}
  a^2 + b^2 &= c^2,
\end{align}
with $a$ and $b$ the catheti, and $c$ the hypotenuse of a right-angled
triangle. \DO{Remember to write the equation in such a way that it is
  incorporated in your sentence. A good rule of thumb is that your
  sentence should make sense when read out loud. Also remember to
  describe the unknowns of the equation.}

\DO{When you want to refer to your equation, you need to reference
  properly which in this case is simply ``see
  \eqref{eqn:pythagoras}'', or \emph{rarely}, ``see
  equation~\eqref{eqn:pythagoras}''. Prefer the first shorter
  version.}

\DO{Remember who your reader is going to be. Typically this is a
  another student in the same course and with similar background to
  you before you started the project. So write in such a way that the
  reader can understand what you have done. This means including
  enough detail so others can understand both the background you have
  read about and what you have done.}